% Part: first-order-logic
% Chapter: introduction
% Section: substitution

\documentclass[../../../include/open-logic-section]{subfiles}

\begin{document}

\olfileid{fol}{int}{sub}

\olsection{Substitutie}

We bespreken een voorbeeld om te laten zien hoe de onderdelen
samenhangen en hoe de ontwikkeling van syntaxis en semantiek de basis
legt voor ons latere, verdergaande onderzoek. Onze
!!{derivation}ssystemen moeten ons in staat stellen de formule
$\Atom{\Obj P}{\Obj a}$ uit $\lforall[\Obj
v_0][\Atom{\Obj P}]{\Obj v_0}$ te !!{derive}. Misschien willen we dit zelfs als
afleidingsregel formuleren. Daarvoor moeten we de regel echter in de
grootst mogelijke algemeenheid kunnen formuleren: niet alleen voor
$\Obj P$, $\Obj a$ en $\Obj v_0$, maar voor iedere !!{formula}~$!A$,
term~$t$ en !!{variable}~$x$. (Bedenk dat !!{constant}s termen zijn;
we zullen echter ook ingewikkeldere termen beschouwen die uit
!!{constant}s en !!{function}s zijn opgebouwd.) We willen dus iets van
de volgende strekking kunnen zeggen: ``Als men
$\lforall[x][!A(x)]$ kan !!{derive}, mag men~$!A(t)$ afleiden---het
resultaat van het verwijderen van $\lforall[x]$ en het vervangen van
$x$ door~$t$.'' Maar wat betekent ``$x$ door~$t$ vervangen'' precies?
Wat is het verband tussen $!A(x)$ en~$!A(t)$? Werkt dit altijd?

Om dit precies te maken definiëren we de bewerking
\emph{substitutie}. Substitutie is in feite lastig: we kunnen niet
zonder meer alle voorkomens van~$x$ in~$!A$ door~$t$ vervangen, en niet
iedere~$t$ mag voor om het even welke~$x$ worden gesubstitueerd. Ook dit
behandelen we met inductieve definities. Zodra dat is gebeurd, wordt de
afleidingsregel ``leid $!A(t)$ af uit $\lforall[x][!A(x)]$'' een
precieze definitie. Bovendien kunnen we aantonen dat dit een goede
afleidingsregel is, in die zin dat $\lforall[x][!A(x)]$ semantisch
$!A(t)$ tot gevolg heeft. Om dat te bewijzen moeten we wederom iets
aantonen waarbij men zich op het eerste gezicht kan afvragen waarom het
nodig is. Dat $\lforall[x][!A(x)]$ semantisch $!A(t)$ tot gevolg heeft,
berust op het feit dat de vraag of $\Sat{M}{!A(t)}$ geldt uitsluitend
afhangt van de waarde van de term~$t$. Als $m$ het !!{element}
van~$\Domain{M}$ is dat door~$t$ wordt aangeduid, dan geldt dus
$\Sat{M}{!A(t)}[s]$ dan en slechts dan als
$\Sat{M}{!A(x)}[\Subst{s}{m}{x}]$. Dit geldt zelfs wanneer $t$
!!{variable}s bevat, maar we moeten het resultaat zorgvuldig
formuleren.

\end{document}
